\section{Origins of the Project}\label{sec:origin}

The original concept for \emph{Il Portale} was filed in January~2013 as a
design specification under the name of Fabrizio Terzi and submitted to the
Free Technology Guild (FTG) journal \emph{Peeragogy in Action} in January~2016
\cite{terzi_peeragogy_2016}.  The 2016 submission (Volume~XII, No.~1)
describes the project as a ``model of use'' combining music education,
peeragogical methodology, and digital media within a compact physical
structure.

The 2013 filing identifies the location already as ``Parco Comunale dei
Caprotti, Bergamo, Via~Tasso''---the same park that anchors the present-day
spatial exploration (see the project dossier in {\tt docs/caprotti/}).

The original document is not a formal academic paper but a project proposal
cast in the FTG format: a ``seed project'' description with technical
drawings, a component breakdown, and a status section that frames the venture
as a working group within the FTG.  The abstract of the 2016 journal issue
contextualises the submission inside the broader peeragogical discourse,
describing the journal as an open editorial space for contributions on
co-generation, recursion theory, and systems thinking in learning.

\paragraph*{Status of the source.}
The 2016 PDF is the only known published version of the original proposal.
No subsequent revision, peer review, or rejection record has been located.
The present document treats it as a historical artefact: it is not corrected,
rewritten, or reinterpreted freely.  Where the contemporary work extends
beyond the 2016 record, the distinction is made explicit.